\chapter{Pendulums and Alternating Current}

\section{Pendulums}\index{pendulum}

A grandfather clock keeps time with a swinging weight, and a playground
swing goes back and forth at its own steady rhythm no matter how high you
pump it. Both are pendulums, and in this chapter you will find out what
sets their rhythm, then see the same mathematics in the electricity in your
walls.

You may be surprised to know that pendulums act the exact same way as roller coasters from the circular chapter, in terms of forces. But, similar to springs,
they oscillate back and forth unless acted upon by a damping force. A pendulum has
an equilibrium position $\theta=0$, hanging straight down. A pendulum also has a maximum angle $\theta_{max}$ from its equilibrium position. It will swing through a set arc in a repetitive motion, which is a
type of motion we call simple harmonic motion, similar to the springs in the oscillations chapter\index{SHM}\index{simple
harmonic motion}. The tension force is always directed towards the pivot. This is
diagramed below in Figure~\ref{fig:pend_dia}
\begin{figure}[H]
	\centering
	\begin{center}
		\begin{tikzpicture}[scale=2]
			\def\L{4} % rod length
			\def\maxAngle{35} \def\bobRad{0.1} % mass-ed bob radius
			\useasboundingbox (-2,0) rectangle (2,4); % manual bounding box

			\coordinate (pivot) at (0,4);
			\coordinate (eqPos) at (0,0);
			\coordinate (midPlus) at ($(pivot)+({-90+\maxAngle/2}:\L)$);
			\coordinate (midMinus) at ($(pivot)+({-90-\maxAngle/2}:\L)$);
			\coordinate (maxPlus) at ($(pivot)+({-90+\maxAngle}:\L)$);
			\coordinate (maxMinus) at ($(pivot)+({-90-\maxAngle}:\L)$);

			\draw[very thick, black!80]
				(pivot) -- (eqPos)
				node[midway, left] {Massless Rod};

			\begin{scope}[rotate around={\maxAngle:(pivot)}]
				\filldraw[black!65] (0,0) circle (\bobRad);
				\draw[very thick, black!60, dashed] (pivot) -- (0,0);
			\end{scope}
			\begin{scope}[rotate around={-\maxAngle:(pivot)}]
				\filldraw[black!65] (0,0) circle (\bobRad);
				\draw[very thick, black!60, dotted] (pivot) -- (0,0);
			\end{scope}

			\filldraw[blue!65]
				(pivot)
				circle
				(0.02)
				node[above, font=\footnotesize] {Pivot};
			\filldraw[black!65]
				(eqPos)
				circle
				(\bobRad)
				node[below, font=\footnotesize] {Equilibrium Position};

			\coordinate (arcStart) at ($(pivot)+({-90+\maxAngle}:\L)$);
			\path (arcStart)
				arc[start angle={-90-\maxAngle}, end angle={-90+\maxAngle}, radius=\L]
				coordinate
				(arcEnd);
			\draw[thick, dashed, red!50, <->]
				(pivot)
				++({-90-\maxAngle}:\L)
				arc[start angle={-90-\maxAngle}, end angle={-90+\maxAngle}, radius=\L]
				node[pos=0.25, below, sloped, font=\footnotesize] {Arc Path};

			\pic [draw, "$\theta $", angle eccentricity=1.5, angle radius=1.2cm, thick, dashed,
			<->] {angle = eqPos--pivot--arcStart};
		\end{tikzpicture}
	\end{center}
	\caption{The anatomy of a pendulum.}
	\label{fig:pend_dia}
\end{figure}
When displaced by an angle $\theta $from the vertical equilibrium position ($\theta
= 0$), gravity provides a \emph{restoring force}that acts along the arc of
motion. This restoring force is given by $F = -mg\sin\theta $, where the negative
sign indicates that the force always points back toward the equilibrium
position. The tension in the string keeps the bob moving along its curved path,
while gravity continually tries to return it to rest. The net force shifts such
that the acceleration arrow is \emph{not}always pointing to the center. The
net force is a combination of the tension (centripetal) force and the
gravitational (tangential), so rather it points at different points along the arc
of the circle, as it is always changing due to the height of the pendulum. This
is demonstrated in Figure~\ref{fig:pend_dia2}. Note the following differences in
energy and forces at the equilibrium position versus the maximum angle positions,
summarized in Table~\ref{tab:pendulum_states}.
\begin{table}[h!]
	\centering
	\begin{tabular}{|c|c|}
		\hline
		\textbf{At $\theta = 0$}                    & \textbf{At $\theta = \theta_{\max}$} \\
		\hline
		$F_{\text{restoring}}= 0$                   & $F_{\text{restoring}}= \text{max}$   \\
		$a_{t}\text{ (tangential acceleration)}= 0$ & $a_{t}= \text{max}$                  \\
		$\text{PE}= 0$                              & $\text{PE}= mgh$                     \\
		$\text{KE}= \text{max}$                     & $\text{KE}= 0$                       \\
		$V = \text{max}$                            & $v = 0$                              \\
		\hline
	\end{tabular}
	\caption{Comparison of pendulum quantities at equilibrium and maximum
	displacement}
	\label{tab:pendulum_states}
\end{table}
Note here that the magnitude of the restoring force $-mg\sin(\theta)$, which is not \emph{directly} proportional to the angular displacement $\theta $. However, for
small angles (less than about $15^\circ$, or 0.26 radians), we can use
the \index{small-angle approximation}small-angle approximation $\sin(\theta) \approx \theta $(in radians), and state
the restoring force as $F \approx -mg\theta $. This allows us to treat the pendulum's motion as simple harmonic motion for small oscillations, where the restoring force
is approximately proportional to the angular displacement.
\begin{figure}[H]
	\centering
	\begin{center}
		\begin{tikzpicture}[scale=2]
			\def\L{4} % rod length
			\def\maxAngle{35} \def\bobRad{0.1} % mass-ed bob radius
			\useasboundingbox (-2,-1) rectangle (2,2); % manual bounding box
			\def\Flen{0.8} % length scale for force arrows

			\coordinate (pivot) at (0,4);
			\coordinate (eqPos) at (0,0);
			\coordinate (midPlus) at ($(pivot)+({-90+\maxAngle/2}:\L)$);
			\coordinate (midMinus) at ($(pivot)+({-90-\maxAngle/2}:\L)$);
			\coordinate (maxPlus) at ($(pivot)+({-90+\maxAngle}:\L)$);
			\coordinate (maxMinus) at ($(pivot)+({-90-\maxAngle}:\L)$);

			\filldraw[black!50] (maxPlus) circle (\bobRad);
			\filldraw[black!50] (maxMinus) circle (\bobRad);

			\filldraw[black!50] (midPlus) circle (\bobRad);
			\filldraw[black!50] (midMinus) circle (\bobRad);

			\filldraw[black!50] (eqPos) circle (\bobRad);

			\coordinate (arcStart) at ($(pivot)+({-90+\maxAngle}:\L)$);
			\path (arcStart)
				arc[start angle={-90-\maxAngle}, end angle={-90+\maxAngle}, radius=\L]
				coordinate
				(arcEnd);
			\draw[thick, dashed, red!50, <->]
				(pivot)
				++({-90-\maxAngle}:\L)
				arc[start angle={-90-\maxAngle}, end angle={-90+\maxAngle}, radius=\L];

			\foreach \P in {maxPlus,maxMinus,midPlus,midMinus,eqPos}{ \draw[->, thick, blue!70] (\P) -- ++(270:\Flen); }
			\node[blue!70, font=\scriptsize, below right]
				at
				($(maxPlus)+(270:\Flen)$)
				{$ m\vec g $};

			\foreach \P in {maxPlus,maxMinus,midPlus,midMinus,eqPos}{ \draw[->, thick, green!60!black] (\P) -- ($(\P)!0.35!(pivot)$); }
			\node[green!60!black, font=\scriptsize, above]
				at
				($(midPlus)!0.35!(pivot)$)
				{$\vec T $};

			\pgfmathsetmacro{\angMaxPlus}{-90+\maxAngle} \pgfmathsetmacro{\angMidPlus}{-90+\maxAngle/2}
			\pgfmathsetmacro{\angMaxMinus}{-90-\maxAngle} \pgfmathsetmacro{\angMidMinus}{-90-\maxAngle/2}

			\draw[->, thick, black!70] (maxPlus) -- ++({\angMaxPlus-90}:0.7*\Flen);
			\draw[->, thick, black!70] (midPlus) -- ++({\angMidPlus-90}:0.7*\Flen);

			\draw[->, thick, black!70] (maxMinus) -- ++({\angMaxMinus+90}:0.7*\Flen);
			\draw[->, thick, black!70] (midMinus) -- ++({\angMidMinus+90}:0.7*\Flen);

			\node[black!70, font=\scriptsize, left]
				at
				($(midPlus)+({\angMidPlus-90}:0.7*\Flen)$)
				{$\vec F_{\text{net}}$};
		\end{tikzpicture}
	\end{center}
	\caption{The net force at different locations of the arc.}
	\label{fig:pend_dia2}
\end{figure}
\begin{Exercise}
	[label={ex:pendulumperiod}, title={Pendulum Period Dependence}] A simple
	pendulum consists of a mass $m$ attached to a massless rod $L $, swinging under the
	influence of gravity. Derive the equation for the period of a simple pendulum. Which
	of the following changes would affect the period $T$ of the pendulum's oscillation?
	\begin{enumerate}[]
		\item Increasing the mass $m$ of the pendulum bob.

		\item Increasing the length $L$ of the string.

		\item Increasing the amplitude (maximum angle) of the swing.

		\item Decreasing the acceleration due to gravity $g $.
	\end{enumerate}
\end{Exercise}
\begin{Answer}
	[ref={ex:pendulumperiod}] The restoring force along the arc is $F = -mg\sin\theta$, so the
	tangential acceleration is $a = -g\sin\theta$. For small angles,
	$\sin\theta \approx \theta$, so $a \approx -g\theta$. The distance along
	the arc from the bottom is $x = L\theta$, so $\theta = \frac{x}{L}$ and
	\begin{equation}
		a \approx -\frac{g}{L}x \label{eq:linearpendulum}
	\end{equation}
	This is the signature of simple harmonic motion, $a = -\omega^2 x$, from
	the oscillations chapter, with
	\begin{equation}
		\omega = \sqrt{\frac{g}{L}}
	\end{equation}
	Therefore, the period $T$ of the pendulum is
	\begin{equation}
		T = \frac{2\pi}{\omega} = 2\pi \sqrt{\frac{L}{g}}
	\end{equation}
	So only changes (b) and (d) affect the period: a longer string or weaker
	gravity makes it longer. The mass cancels out, and for small swings the
	amplitude does not matter.
\end{Answer}
The period of a simple pendulum depends on its length and on gravity,
but not on the mass of the bob or (for small swings) the amplitude. Galileo
noticed this around 1600, supposedly by timing a swinging lamp in a
cathedral against his pulse, and it is why pendulums made the first
accurate clocks. Like a mass on a spring, the pendulum's angle follows a
sine wave: $\theta(t) = \theta_{max}\sin\left(\sqrt{\frac{g}{L}}\,t + \phi_{0}\right)$.

\begin{Exercise}
	[label={ex:pendulumMoon}, title={Pendulum on the Moon}] A simple pendulum has
	a period of 2 seconds when measured on Earth. If the same pendulum were taken to
	the Moon, where the acceleration due to gravity is approximately $\frac{1}{6}$ of Earth's,
	what would be its new period?
	\begin{enumerate}[label=\alph{*}]
		\item Approximately 0.8 seconds

		\item Approximately 2 seconds

		\item Approximately 4.9 seconds

		\item Approximately 6.8 seconds
	\end{enumerate}
\end{Exercise}
\begin{Answer}
	[ref={ex:pendulumMoon}] The correct answer is (c): approximately 4.9 seconds.

	The period $T$ of a simple pendulum is given by the formula:
	\[
		T = 2\pi \sqrt{\frac{L}{g}}
	\]
	where $L$ is the length of the pendulum and $g$ is the acceleration due to gravity.

	On Earth, the period is 2 seconds, so we can set up the equation:
	\[
		2 = 2\pi \sqrt{\frac{L}{g_{\text{Earth}}}}
	\]
	This shows that $T$ is inversely proportional to the square root of $g $. We can
	rearrange this equation to solve for $L $. So if $g$ decreases by a factor of 6,
	then $T$ will increase by a factor of $\sqrt{6}$.

	Now, on the Moon, the acceleration due to gravity is approximately $\frac{1}{6}
	g_{\text{Earth}}$. We can substitute this into the period formula:
	\[
		T_{\text{Moon}}= 2\pi \sqrt{\frac{L}{g_{\text{Moon}}}}= 2\pi \sqrt{\frac{L}{\frac{1}{6}g_{\text{Earth}}}}
		= 2\pi \sqrt{\frac{6L}{g_{\text{Earth}}}}= \sqrt{6}\cdot 2\pi \sqrt{\frac{L}{g_{\text{Earth}}}}
		= \sqrt{6}\cdot T_{\text{Earth}}
	\]
	Since our original period on Earth is 2 seconds, we have:
	\[
		T_{\text{Moon}}= \sqrt{6}\cdot 2 \approx 4.9 \text{ seconds}
	\]
	Which is equal to option (c).
\end{Answer}

\section{Alternating Current}\index{alternating current}
We have talked about Alternating Current (AC) in previous chapters, but now we will
explore its connection to oscillations. AC is an electric current that
periodically reverses direction, in contrast to Direct Current (DC), which flows
in a single direction and provides constant voltage. The voltage and current in an
AC circuit vary sinusoidally with time, making them ideal for modeling using
oscillatory functions. 

In a DC circuit, the voltage is constant: $V = V_0$.

In an AC circuit, the voltage oscillates exactly like the position of a mass on a spring:

\begin{equation}
V(t) = V_{peak} \sin(\omega t + \phi)
\end{equation}

where
\begin{itemize}
	\item $V_{peak}$ is the maximum voltage reached in either direction
	\item $\omega$ is the angular frequency, related to standard frequency by $\omega ={\frac {2\pi }{T}}={2\pi f}$
	\item $\phi$ which is the phase constant, or the starting position of the wave
\end{itemize}
Note that, because of the trigonometric nature of sine, the wave fluctuates between $V_{peak}$ and $-V_{peak}$. 

%https://en.wikipedia.org/wiki/Alternating_current 
As you saw in the alternating current chapter, the average power from a
sine-wave voltage is the same as from a steady voltage of
$V_{rms} = \frac{V_{peak}}{\sqrt{2}}$.
\subsection{Impedance}\index{impedance}
In DC, we use Resistance ($R$). In AC oscillations, we use \textit{Impedance}. Impedance is the ``total opposition'' to the oscillation and is a combination of standard resistance and Reactance ($X$).

Capacitors oppose changes in voltage, which we call Capacitive Reactance:$$X_C = \frac{1}{\omega C}$$

Inductive Reactance comes from inductors which oppose changes in current (acting like ``electrical inertia'').$$X_L = \omega L$$

Similarly, in AC circuits, the voltage and current waves do not always peak at the same time. 
\begin{itemize}
	\item In a purely inductive circuit, the current falls behind the voltage by $90^\circ$ ($\pi/2$). 
	\item In a purely capacitive circuit, the current leads the voltage by $90^\circ$ ($\pi/2$).
\end{itemize}
\subsection{Comparing Mass-Springs and AC}
\begin{table}[!ht]
    \centering
    \begin{tabular}{|l|l|}
    \hline
        Mechanical System (Mass-Spring) & Electrical System (AC Circuit) \\ \hline
        Displacement (x) & Charge (q) \\ \hline
        Velocity (v) & Current (i) \\ \hline
        Mass (m) & Inductance (L) \\ \hline
        Spring Constant (k) & Reciprocal Capacitance (1/C) \\ \hline
        Damping/Friction (b) & Resistance (R) \\ \hline
    \end{tabular}
\end{table}
\section{Summary}
In this chapter, we have discussed oscillations via Simple Harmonic Motion, mass-springs, spring energy, pendulums, and AC currents.
\begin{itemize}
	\item \textbf{Simple Harmonic Motion (SHM)}: periodic motion in which a restoring force, proportional to displacement, pulls an object back toward equilibrium. One full cycle takes $T = 2\pi/\omega$.
	
	\item \textbf{Hooke's Law}: For an ideal spring, $F_s = -kx $. This linear relationship leads to an angular frequency of $\omega = \sqrt{k/m}$.
	\item \textbf{Pendulums}: For small swings, $T = 2\pi\sqrt{L/g}$, independent of mass and amplitude.
	\item \textbf{Energy Dynamics}: In a closed system, energy oscillates between Kinetic ($K$) and Potential ($U$). Total mechanical energy remains constant ($E = \frac{1}{2}kA^2$), with $K$ peaking at equilibrium and $U$ peaking at maximum displacement.
	\item \textbf{AC circuits}: Voltage oscillates like a mass on a spring, with inductance playing the role of mass and $1/C$ the role of the spring constant.
\end{itemize}

This chapter leads us into our next chapter on orbits, the oscillations of planetary motion.

\textbf{Worked example.} What is the period of a pendulum 1 m long?

\textbf{Solution.} $T = 2\pi\sqrt{\frac{1}{9.8}} \approx 2.01$ s: about one
second per swing in each direction.

\begin{Exercise}[title={Shorter Pendulum}, label=pend3]
A pendulum has length $0.25$ m. What is its period?
\end{Exercise}
\begin{Answer}[ref=pend3]
$T = 2\pi\sqrt{\frac{0.25}{9.8}} \approx 1.00$ s.
\end{Answer}
